%=======================================================================
%  Paper V — Section 1 (Abstract + Introduction)   revised
%  Lines marked  % VERIFY  contain values to confirm against your tables
%=======================================================================
%\documentclass[preprint,12pt]{elsarticle}

%\usepackage[T1]{fontenc}
%\usepackage[utf8]{inputenc}
%\usepackage{amsmath,amssymb}
%\usepackage{graphicx}
%\usepackage{booktabs}
%\usepackage{natbib}
%\usepackage{siunitx}
%\sisetup{per-mode=symbol, range-phrase={--}, range-units=single}
%\usepackage[hidelinks]{hyperref}

%------------------------- notation macros -----------------------------
%\newcommand{\potvarn}{\tilde{\sigma}_{U}}
%\newcommand{\denshead}{\varrho_{\mathrm{head}}}
%\newcommand{\densrest}{\varrho_{\mathrm{rest}}}
%\newcommand{\densR}{\varrho_{R}}
%\newcommand{\densbulk}{\varrho_{\mathrm{bulk}}}
%\newcommand{\Ddens}{\Delta\varrho}
%\newcommand{\Dmass}{\Delta M}
%\newcommand{\VR}{V_{R}}
%\newcommand{\Vtot}{V_{\mathrm{tot}}}
%\newcommand{\Mtot}{M_{\mathrm{tot}}}
%\newcommand{\ampfac}{\mathcal{A}}
%\newcommand{\Ufull}{U^{\mathrm{grav}}_{\mathrm{full}}}
\newcommand{\Ureg}{U^{\mathrm{grav}}_{R}}
\newcommand{\Urot}{U^{\mathrm{rot}}}

%\journal{Icarus}

%\begin{document}

%\begin{frontmatter}

%\title{What potential-dispersion minimization constrains in the interior of a small body}

%\author{...}   % add author block

%=======================================================================
%  ABSTRACT  --  target 250--300 words (this version: 296)
%=======================================================================
{\noindent\large\bfseries Abstract}\\[0.4em]
{\noindent
\citet{richardson2014} argued that on a small body whose surface has
relaxed under downslope regolith flow the bulk density can be estimated
by minimizing the dispersion of the surface geopotential, and
\citet{kanamaru2019} extended the argument from a single mean density to
a two-region interior with the total mass held fixed. Paper~IV
established the property that makes such a minimization meaningful: the
normalized geopotential dispersion $\potvarn$ is far less sensitive to
mesh resolution than the mean surface slope. We ask what the extension
to interior structure actually constrains.

Regions are built by exact plane clipping, verified against closed-form
solids, a complement closure, a vertex predicate and a manifold criterion,
with a fifth containment test where a shape model has more than one
component, and externally against the published lobe volume of 67P. Of five bodies, only Itokawa gives a dispersion
improvement above the $0.9$--$5.9\%$ mesh-resolution sensitivity
envelope of Paper~IV: $16.9\%$ at the detected neck. Eros, Bennu, 67P
and Arrokoth return $<10^{-3}\%$ to $1.3\%$, at or below that envelope.
At the dividing plane $x=\SI{150}{\metre}$ of \citet{kanamaru2019} we
recover $\denshead=2725$ and $\densrest=\SI{1858}{\kilo\gram\per\cubic\metre}$,
a contrast of $1.47$ against their $1.46$, and a centre-of-mass offset of
$16.5\pm\SI{2.9}{\metre}$, consistent with the $\SI{16}{\metre}$ from
equipotential fitting and with the $21\pm\SI{12}{\metre}$ from YORP spin-up
\citep{lowry2014}---though that determination is itself only $1.75$ standard
deviations from zero, and its interval contains every plane of our sweep.

The mass constraint makes the excess mass $\Dmass=\Ddens\VR$ identically
$\Mtot-\densrest\Vtot$, so a statement about one is a statement about the
other. Over the eleven planes of a thirteen-plane sweep on Itokawa that lie
below the grain-density ceiling, $\denshead$ varies by $\pm13.2\%$ and
$\densrest$ by $\pm0.43\%$. In coefficients of variation the ratio is $32$;
the identity contributes $11.8$ of it and the subtraction that defines
$\Ddens$ removes $2.8$, leaving a factor of $7.5$ that reflects genuine
compensation between contrast and region volume. Including the curvature sensitivity interval
gives $\densrest=1858\pm\SI{27}{\kilo\gram\per\cubic\metre}$.
Eight uniform synthetic contact binaries show the invariance emerges
only as the signal strengthens. Dispersion minimization does not
constrain its own dividing plane, and Itokawa fails the first of the
four relaxation conditions. Conservation of the mass dipole is excluded at twenty-one standard errors,
so the agreement with YORP is an independent check and not a second reading
of the same quantity. Dispersion minimization does not constrain its own
dividing plane: left to itself it runs to a head density standing $11.2$
standard deviations above the LL grain density. Itokawa fails the first of
the four relaxation conditions, but that condition governs the conditioning
of a variable-mass inversion and has no counterpart here, and a control run
with the rotational term removed confirms that the fixed-mass objective
retains an interior minimum. On Arrokoth, which has no measured mass, the
sense of the recovered contrast reverses within the published range of bulk
density, between $235$ and $\SI{250}{\kilo\gram\per\cubic\metre}$.
%\end{abstract}


%-----------------------------------------------------------------------

}

\medskip
\noindent\textbf{Keywords:} 
asteroids,
interior structure,
small bodies,
surface geopotential,
polyhedral gravity,
inverse problems,
identifiability,
shape models,
25143 Itokawa,
486958 Arrokoth,
67P/Churyumov-Gerasimenko,
101955 Bennu,
433 Eros

%\end{frontmatter}

%=======================================================================
\section{Introduction}
\label{sec:intro}

The interior of a small body is not directly observable. Where a
spacecraft has tracked the body, radio science fixes the total mass and,
on the most favourable targets, the low-degree gravity field
\citep{yeomans2000,patzold2016,scheeres2020}; but the mass alone is one
number, and the gravity coefficients that would resolve internal
structure are degenerate with the shape unless the spacecraft flies very
close for a long time. Where no spacecraft has visited, even the mass is
absent. Indirect routes are therefore of interest, and one of them uses
the surface itself.

%-----------------------------------------------------------------------
\subsection{Relaxation as an inference tool}
\label{sec:intro:relaxation}

\citet{richardson2014} set out the argument in full. On a rotating
irregular body covered by mobile regolith, downslope flow triggered by
seismic shaking or volatile activity acts as a negative-feedback system,
eroding local topography toward a flat equipotential surface. If that
process has run to completion, the surface geopotential should be nearly
uniform, and the bulk density can be estimated by finding the value that
minimizes its dispersion. They state four conditions under which the
argument holds: the mean rotational force must be a significant fraction
of the mean gravitational force; a sufficiently thick, low-cohesion,
mobile regolith layer must exist over most of the surface; a downslope
flow disturbance source must be present; and sufficient time must have
passed since the last major surface alteration. They are explicit that
the result is neither a direct nor an indirect density measurement, but
the density corresponding to the most eroded state of the body given its
present shape, topography and spin. We return in
Section~\ref{sec:discussion} to the standing of these four conditions on
the bodies examined here, and note at the outset that the single body
that responds to the inversion, Itokawa, does not satisfy the first of
them.

\citet{kanamaru2019} extended the argument from a single mean density to
an interior distribution. Holding the total mass at its measured value
and dividing Itokawa into a head and a body at $x=\SI{150}{\metre}$,
they minimized the area-weighted potential variance over the global
surface with respect to the one remaining free parameter, and recovered
a head density of $2730$ and a body density of
$\SI{1870}{\kilo\gram\per\cubic\metre}$, a contrast of $1.46$. A
companion study using the smooth terrains rather than the whole surface
\citep{kanamaru2019pss} obtained $2450$ and
$\SI{1930}{\kilo\gram\per\cubic\metre}$, a contrast of $1.27$, and,
importantly, found that a model in which the two lobes have comparable
densities separated by a compressed neck fits marginally better than one
in which the lobes differ, and that a three-region model treating head,
neck and body separately is largely unconstrained; only when the
independently estimated centre-of-mass offset from YORP spin-up is
imposed do their models require the head to exceed the global average.
\citet{motooka2012} had earlier attempted to fit the global shape of
Itokawa to a single equipotential surface. An independent line of
evidence comes from \citet{lowry2014}, whose measured YORP spin-up
requires a centre-of-mass displacement of
$\SI{21\pm12}{\metre}$ along the long axis. Because this is the only
constraint on Itokawa's interior not derived from an equipotential
argument, it is the only genuinely external check available to us, and
we treat it as such in Section~\ref{sec:itokawa}.

Paper~IV of this series supplied a property on which any such
minimization depends. Comparing meshes of the same body decimated from
roughly $42{,}000$ to $2{,}000$ facets, the area-weighted mean slope
varied by $1.8\%$ to $25.6\%$ across the test set while the normalized
geopotential dispersion varied by only $0.9\%$ to $5.9\%$, reaching a
plateau between the two finest meshes. The endpoints of these two ranges
are attained on different bodies and should not be divided against one
another; the per-body ratios are given in Paper~IV. A quantity that
changes with the arbitrary resolution of the mesh cannot support an
inference; the dispersion, over the tested range, does not.

Section~\ref{sec:discussion:ellipsoid} tests that comparison against an
analytic reference and finds it, if anything, understated. We refer to
the $0.9$--$5.9\%$ range throughout as the \emph{mesh-resolution
sensitivity envelope}. It measures how far $\potvarn$ moves when the
discretization changes; it is neither a statistical noise level nor a
detection threshold, and we use it only as a scale against which an
improvement may be judged large or small. Constructing a proper null
distribution would require either a noise model for the shape data or a
placebo ensemble of randomly placed dividing surfaces; neither is
attempted here, and the absence is a limitation we state plainly rather
than a gap we conceal.

%-----------------------------------------------------------------------
\subsection{The question this paper asks}
\label{sec:intro:question}

That the extension works on Itokawa has been reported. What it
constrains has not been examined. The dividing plane is chosen by the
investigator, and no published study varies it and reports how the
recovered density responds. The question matters because the density of
a region is not an observable in the way the total mass is: it is the
ratio of an inferred excess mass to a volume that the investigator has
defined.

The point can be made before any computation. With $\Mtot$ and $\Vtot$
held fixed, conservation of mass gives
%
\begin{equation}
\Mtot \;=\; \densR\VR + \densrest\!\left(\Vtot-\VR\right)
      \;=\; \densrest\Vtot + \Ddens\,\VR ,
\label{eq:massbalance}
\end{equation}
%
with $\Ddens\equiv\densR-\densrest$, so that the excess mass
$\Dmass\equiv\Ddens\VR$ satisfies the exact identity
%
\begin{equation}
\Dmass \;=\; \Mtot-\densrest\Vtot
       \;=\; \left(\densbulk-\densrest\right)\Vtot ,
\qquad \densbulk\equiv\Mtot/\Vtot .
\label{eq:excessidentity}
\end{equation}
%
Equation~\eqref{eq:excessidentity} involves no property of the objective
function. It says that $\Dmass$ and $\densrest$ are related by a
bijective affine map, that a two-region model with fixed total mass has
exactly one free parameter, and that the pair $(\densR,\VR)$ is
degenerate along the hyperbola $\Ddens\VR=\mathrm{const}$. Any claim
that one of these quantities is ``more stable'' than another must
therefore be accompanied by the amplification factor relating them,
%
\begin{equation}
\ampfac \;\equiv\; \frac{\densrest\Vtot}{\Dmass}
        \;=\; \frac{\densrest}{\densbulk-\densrest}
        \;=\; \frac{1}{f\,(C-1)},
\qquad f\equiv\VR/\Vtot,\quad C\equiv\densR/\densrest ,
\label{eq:amplification}
\end{equation}
%
which on Itokawa at the detected neck takes the value
$\ampfac=12.0$.
We show in Section~\ref{sec:constrains} that reporting the relative
scatter of $\densrest$ without $\ampfac$ overstates the precision of the
inversion by exactly this factor, and that $\ampfac$ itself is amplified
by $\densbulk/(\densbulk-\densrest)\simeq13$ relative to errors in
$\densrest$, so that a $0.13\%$ shift in the background density moves
$\ampfac$ by $1.7\%$.

We therefore treat the method as an object of study rather than a tool.
Five bodies are analysed---Itokawa, Eros, Bennu, 67P and
Arrokoth---together with eight synthetic contact binaries constructed so
that the strength of the signal can be varied while everything else is
controlled: uniform in density by construction, volumes within $7\%$ of
one another, and every plane sweep spanning the same factor $1.90$ in
region volume. The synthetic experiment is necessary because no further
real body can settle the question: we cannot choose in advance whether
it will carry a signal, and no body's true interior is known, so
agreement would establish only that something had happened twice. We
note the corresponding limitation: because the synthetic bodies are
uniform, they bound the false-positive rate but do not measure recovery
of a known contrast, and the detection limit of the method is therefore
not established here.

%-----------------------------------------------------------------------
\subsection{A construction error and how it was found}
\label{sec:intro:construction}

One result of this work is negative and concerns the region construction
itself, and we state it before the contributions because it conditions
all of them.

An earlier implementation closed each region by fanning its boundary
edges to an interior point. That construction is watertight, has an
exact tetrahedral volume, and passes a volume-closure test---and it
encloses the wrong solid. Fanning to an interior point produces a radial
cone over the surface patch, not the plane-cut lobe intended. On a unit
sphere cut at $x=0.5$ the cone is a spherical sector of volume $1.0472$
and centroid $0.5625$, against $0.6545$ and $0.6750$ for the spherical
cap; on Itokawa the resulting error reversed the sign of the recovered
contrast. The volume-closure test then in use compared the fan mesh
against the same tetrahedra and could not detect the discrepancy. Only a
\emph{predicate} test---requiring every vertex of the region to satisfy
its defining inequality---distinguishes the two constructions.

The present implementation builds regions by exact plane clipping and is
verified in four independent geometric layers---against closed-form solids,
by complement closure, by the vertex predicate and by a manifold
criterion---with a fifth containment test where a shape model has more than
one component, an arithmetic audit of the mass balance on every solution,
and externally against the published lobe volume of 67P, where our head region holds $26.6\%$ of
the body against the $27\%$ given by \citet{jorda2016}. We record in
Section~\ref{sec:methods} that this external agreement is weaker than it
appears, because a single plane assigns the entire neck to one side
whereas the lobe fit of \citet{jorda2016} assigns $66\%$ and $27\%$ of
the volume to the two lobes and leaves the remaining $7\%$, the neck, in
neither; moving the plane across the neck changes our percentage by
$10.7$ points, and we report that sensitivity so that the reader may
judge it.

%-----------------------------------------------------------------------
\subsection{What is new here}
\label{sec:intro:new}

To be plain about the boundary: the relaxation argument and the
potential-variance technique are \citet{richardson2014}; the extension
to a two-region interior is \citet{kanamaru2019}; the polyhedron gravity
is \citet{werner1996}; and the finding that Itokawa's head is the denser
lobe is \citet{lowry2014} and \citet{kanamaru2019}, reproduced here but
not discovered here. \citet{scheeres2020} describe that finding as
tentative and not directly tested, which is part of what motivates an
independent implementation.

The contributions of this paper are the following.

\begin{enumerate}
\item An exact plane-clipping region construction with four verification
layers, of which the predicate and manifold tests can each fail in ways
the other three cannot, together with external validation against the
published lobe volume of 67P and an explicit account of how the neck is
apportioned (Section~\ref{sec:methods}).

\item The demonstration that the density recovered by the inversion is
resolution-stable over the meshes tested---a property distinct from, and
not implied by, the stability of the dispersion established in
Paper~IV---with no change across the three Itokawa meshes that share a
dividing plane and none across the three Arrokoth meshes at a fixed plane.
``No change'' means below the resolution of the scan, which bounds it at
$\SI{12.5}{\kilo\gram\per\cubic\metre}$, or $0.45\%$, on Itokawa; on
Arrokoth the recovered density is identical on all three meshes and the
complementary density moves by $\SI{0.4}{\kilo\gram\per\cubic\metre}$, or
$0.1\%$ (Sections~\ref{sec:itokawa} and~\ref{sec:nulls}).

\item The decomposition of the apparent stability of the inversion into
its algebraic and physical parts. Equation~\eqref{eq:excessidentity}
fixes a factor of $11.8$ between the coefficients of variation of $\Dmass$
and of $\densrest$, and the definition of $\Ddens$ as a difference inflates
its coefficient of variation over that of $\denshead$ by a further factor of
$2.8$; only the factor of $7.5$ by which $\Dmass$ is steadier than $\Ddens$
reflects compensation between contrast and region volume and is a property
of the objective function. The three combine exactly to the net factor of
$32$ between $\denshead$ and $\densrest$ over the eleven planes below the
grain-density ceiling (Section~\ref{sec:constrains}).

\item The identification of the background density, and with it the mass
anomaly to which \eqref{eq:excessidentity} ties it exactly, as the
quantity the method constrains in the strong-signal regime, together
with the demonstration by controlled synthetic experiment that this
holds in that regime rather than of the method as such
(Section~\ref{sec:constrains}).

\item The exclusion of the competing hypothesis that the conserved
quantity is the mass dipole $D=\Dmass\,(x_{R}-x_{cf})$. Dipole
conservation requires $\mathrm{d}\log\Dmass/\mathrm{d}\log\VR>0$ and
hence a scaling slope $s>-1$; quantitatively it requires $s=-0.756$ over
the eleven Itokawa planes below the grain-density ceiling, and the measured
$s=-1.141\pm0.018$ lies on the opposite side of $-1$, twenty-one standard
errors away (Sections~\ref{sec:itokawa} and~\ref{sec:constrains}).

\item The result that dispersion minimization does not constrain its own
dividing plane---at the outermost plane of the sweep the head density it
prefers stands $11.2$ standard deviations above the LL grain density---and
that an external physical bound is therefore required
(Section~\ref{sec:plane}).

\item Four weak-signal cases, establishing that the method does not
manufacture structure when the plane is correctly placed, and one worked
case showing what a misplaced plane costs (Section~\ref{sec:nulls}).

\item The finding that on a body without a measured mass the assumed
bulk density places the inferred contrast at the threshold of a sign
reversal, and the suggestion that sensitivity to that assumption may
serve as a secondary diagnostic (Sections~\ref{sec:nulls}
and~\ref{sec:constrains}).

\item An independent test of the Paper~IV resolution result against a
homogeneous ellipsoid, whose surface gravity has a closed form, removing
the need to treat the finest mesh as the truth
(Section~\ref{sec:discussion:ellipsoid}).
\end{enumerate}

Section~\ref{sec:methods} sets out the density model, the region
construction and its verification. Section~\ref{sec:itokawa} presents
the single strong signal and its stability properties.
Section~\ref{sec:nulls} presents the four weak-signal cases.
Section~\ref{sec:constrains} is the core of the paper: what the method
constrains, and the synthetic experiment that establishes in which
regime it does so. Section~\ref{sec:plane} shows that the dividing plane
is not self-determined. Section~\ref{sec:failures} catalogues the
failure modes. Section~\ref{sec:discussion} discusses the standing of
the relaxation premise on these bodies, Section~\ref{sec:guidelines}
sets out what the results imply for anyone applying the method, and
Section~\ref{sec:conclusions} states the conclusions.

%\end{document}
